\section{評価計画}

\subsection{データセットと split}

Phase 0 では RanSAP/RanSMAP の schema、license、run 単位、label 意味を監査し、
10 秒統計へ落とせることを確認する。Phase 1 は監査済み RanSAP を主データセットにし、
benign run だけで AE を学習する。calibration split も benign から取り、
ransomware label は評価にのみ使う。split は隣接窓の random split を避け、
run、ransomware family、storage condition、OS/BitLocker 条件で分ける。
Phase 2 以降で RanSMAP の storage-only 部分、SNIA IOTTA または UMass の
benign trace を追加し、false positive と workload shift を評価する。

\begin{table}[tb]
  \centering
  \caption{初期データ利用計画}
  \label{tab:dataset-plan}
  \begin{tabularx}{\linewidth}{lYYYY}
    \toprule
    データ & 攻撃 & benign & entropy & 主用途 \\
    \midrule
    RanSAP & yes & yes & write entropy & 初期 feasibility、entropy upper-bound ablation。 \\
    RanSMAP & yes & yes & storage plus memory context & storage-only 限界分析。 \\
    SNIA IOTTA & no & yes & usually no & false positive、多様な benign workload。 \\
    UMass traces & no & yes & no & OLTP/search などの benign stress。 \\
    NapierOne & no & file corpus & computable offline & 高 entropy benign replay と校正。 \\
    \bottomrule
  \end{tabularx}
\end{table}

\subsection{未知脅威評価}

未知脅威への対応は、任意の将来攻撃を検出できるという強い主張ではなく、
攻撃ラベルに依存しない検出器が、既知 family に過適合した参照分類器とは
異なる失敗様式を持つかを測る proxy claim として扱う。AE の training split と
calibration split は benign run だけから作り、ransomware label、family label、
攻撃開始時刻は threshold tuning に使わない。ransomware run は評価専用に保持し、
family、variant、storage condition ごとに recall、time to detect、alert 前被害量を
層別集計する。

family holdout は二つの意味を分ける。教師なし AE では、held-out family は
学習に使われない攻撃評価 strata であり、AE が攻撃 family label を見ないことを
確認するための集計単位である。教師あり参照モデルでは、既知 family の攻撃 label で
学習・調整し、held-out family で評価する。したがって、AE と教師あり参照の
family holdout 結果は同一条件の勝敗ではなく、未知・変種攻撃を想定したときに
ラベル依存手法と正常逸脱手法がどのように崩れるかを比較する診断である。

\subsection{ベースライン}

AE の価値を評価するには、単純な規則を必ず比較対象にする。baseline は
同じ科学的問いを答えるものだけを同列に置く。

\begin{itemize}
  \item Rule baselines: write ratio threshold、write bytes または write IOPS の
        rolling z-score、平均 LBA/平均転送長 shift、cheap telemetry が
        ある場合の compression/entropy-like threshold。
  \item Unsupervised baselines: Isolation Forest、one-class SVM、PCA または
        robust covariance など、benign training split だけで調整できる手法。
  \item Supervised reference models: 小型 decision tree、random forest、XGBoost 相当。
        これらは label を使える場合の参照上限であり、one-class AE の
        deployment claim と同じ条件の baseline ではない。特に未知 family や
        family holdout 条件では、教師あり分類器の結果は既知ラベル依存の参照値として
        分離する。
\end{itemize}

AE が単に ``busy workload'' を検出しているだけなら、write throughput rule に
負けるはずである。AE が supervised reference に負けても、そのこと自体は
one-class 検出の否定ではないが、ラベル付き storage-level classifier の方が
組み込み経路に適する可能性を示す。比較表では、feature availability、label 使用、
threshold tuning split、tuning budget を明示する。

\subsection{主要指標}

AUROC は補助指標に留める。実運用で重要なのは、低 false positive で早く検出できるか
である。主要指標は以下とする。

\begin{itemize}
  \item fixed false positive budget における recall。
  \item false alarms per volume per day。
  \item attack start から初回 alert までの time to detect。
  \item raw trace と攻撃ラベルから後処理で推定する、alert 前の overwritten bytes
        または affected LBA fraction。
  \item precision-recall curve と AUPRC。
  \item channel-wise reconstruction-error contribution。
  \item MNN 変換後の score parity、weight size、per-volume detector-data size、
        transient scratch。
  \item volume 数 $V$ に対する aggregate memory と 10 秒あたり CPU 時間。
\end{itemize}

ここで false positive は分類上の誤陽性、false alarms per volume per day は
運用上の alert 率として区別する。device-fit のための volume scaling では、
共有 runtime、共有 weight、volume ごとの detector state、transient scratch を分け、
式 \eqref{eq:aggregate-cost} の形で総コストを報告する。$B_{\mathrm{volume}}$ は
500 KB の per-volume detector-data 予算、$Q$ は同時に確保する inference scratch slot 数である。

\begin{equation}
\begin{aligned}
  M_{\mathrm{persistent}}(V) =
  M_{\mathrm{shared}} + M_{\mathrm{w,shared}}
  + V \left(M_{\mathrm{w,volume}} + M_{\mathrm{state}}\right), \\
  M_{\mathrm{peak}}(V,Q) =
  M_{\mathrm{persistent}}(V) + Q M_{\mathrm{scratch}}, \\
  C_{\mathrm{total}}(V) =
  V C_{\mathrm{volume}} / 10\,\mathrm{s}.
\end{aligned}
  \label{eq:aggregate-cost}
\end{equation}

$M_{\mathrm{shared}}$ は共有 runtime/library と共通 heap、$M_{\mathrm{w,shared}}$ は
全 volume で共有できる model weight 表現、$M_{\mathrm{w,volume}}$ は volume ごとに
複製される weight 表現、$M_{\mathrm{state}}$ は retained input statistics、
normalization state、threshold、score history である。weight を共有できる実装では
$M_{\mathrm{w,volume}}=0$ とし、共有できない実装では $M_{\mathrm{w,shared}}=0$ とする。
$M_{\mathrm{scratch}}$ は activation tensor、operator workspace、output tensor などの
reusable inference slot 1 個分である。per-volume gate は、共有 weight なら
$(M_{\mathrm{w,shared}}/V)+M_{\mathrm{state}}\le B_{\mathrm{volume}}$、複製 weight なら
$M_{\mathrm{w,volume}}+M_{\mathrm{state}}\le B_{\mathrm{volume}}$ として報告する。
$C_{\mathrm{volume}}$ は 1 volume の 10 秒 frame 更新と推論に必要な CPU 時間である。
target device の具体的 budget が未定の場合でも、
$V \in \{1,8,32,128,2000\}$ などの sweep を示し、単一 volume の
500 KB detector-data fit だけで device-fit claim を昇格しない。

\subsection{アブレーション}

本研究で最も重要なアブレーションは entropy/compression を外す実験である。
entropy はランサムウェア検出に強い特徴になり得るが、payload access が必要な場合、
ブロック装置内で高コストになり、暗号化済みボリュームや圧縮済み benign file では
識別力が落ちる可能性がある。RanSAP の write entropy を使う結果は、
cheap device telemetry が確認されるまでは upper-bound result として分離する。
したがって、表 \ref{tab:ablation} の ablation を必須にする。

\begin{table}[tb]
  \centering
  \caption{必須アブレーション}
  \label{tab:ablation}
  \begin{tabularx}{\linewidth}{lY}
    \toprule
    実験 & 検証する問い \\
    \midrule
    entropy/compression なし & metadata-only でも検出可能か。 \\
    mean LBA なし & 平均アドレス移動に依存しすぎていないか。 \\
    mean length なし & 平均 I/O サイズの寄与はどれだけか。 \\
    count/bytes intensity 分離 & 忙しい benign workload を誤検出していないか。 \\
    AE family comparison & AE-0/1 MLP、AE-2 two-level Dense、AE-4 Conv1D、AE-3 GRU-context、AE-5 CNN-GRU のどれだけ複雑な architecture が必要か。 \\
    $N=6,12,24$ & $N \times 10$ 秒の文脈長と volume あたり 500 KB detector-data 制約の trade-off。 \\
    volume 数 scaling & per-volume では軽い検出器が全体では重くならないか。 \\
    family holdout & 未知または変種の攻撃に対する教師なし検出価値は残るか。 \\
    quantized MNN vs offline & 量子化と変換で score が崩れないか。 \\
    \bottomrule
  \end{tabularx}
\end{table}
